% claim-ref: sedum-loves-fire#callus-duration
% claim-ref: sedum-loves-fire#callus-preparation
% claim-ref: sedum-loves-fire#failure-cue-no-resistance
% claim-ref: sedum-loves-fire#failure-cue-scorch
% claim-ref: sedum-loves-fire#failure-cue-soft-dark-base
% claim-ref: sedum-loves-fire#insert-cutting
% claim-ref: sedum-loves-fire#leaf-method
% claim-ref: sedum-loves-fire#prepare-tools-and-medium
% claim-ref: sedum-loves-fire#primary-method
% claim-ref: sedum-loves-fire#root-check
% claim-ref: sedum-loves-fire#rooting-medium
% claim-ref: sedum-loves-fire#rooting-moisture
% claim-ref: sedum-loves-fire#rooting-temperature
% claim-ref: sedum-loves-fire#stem-cutting-length
% claim-ref: sedum-loves-fire#strip-lower-stem
% claim-ref: sedum-loves-fire#terminal-cutting-root-initiation
% claim-ref: sedum-loves-fire#transfer-and-acclimate
% claim-ref: sedum-loves-fire#transplant-readiness
\CompanionHeader{PROPAGATION}{Sedum 'Love's Fire'}{Candidate Sedum adolphi 'LOVE'S FIRE'}{Outdoor grow bag; reported label is not a verified cultivar identification.}
\Context{Terminal stem cutting. General guidance; observe parent readiness before acting.}
\begin{minipage}[t]{.625\linewidth}\raggedright\vspace{0pt}
\Step{1}{Prepare}{Clean sharp pruners; fill a small draining pot with cactus potting soil or perlite. This is temporary rooting medium, not the established-plant mix. \textbf{[4]}}
\Step{2}{Select \& cut}{Take a healthy 3-6 in stem tip. This extension range is general succulent guidance, not a cultivar-specific requirement. \textbf{[2,4]}}
\Step{3}{Expose the lower stem}{Remove leaves from the lower half. Retain healthy upper growth and keep the cut clean. \textbf{[4]}}
\Step{4}{Dry \& callus}{Rest dry on a tray or in an empty cup for several days. Proceed by a dry, callused surface, not by date alone. \textbf{[4]}}
\Step{5}{Place upright}{Insert the callused lower stem upright in draining medium. Avoid water-glass rooting; no numeric insertion depth is established. \textbf{[4]}}
\Step{6}{Maintain conditions}{Avoid saturation and standing water. Use slight moisture and water sparingly until roots begin. The patent used 75 F (24 C; F rounded), not a universal home requirement. \textbf{[2,3]}}
\Step{7}{Inspect without uprooting}{Check very gently for anchoring resistance; do not repeatedly lift the cutting. New roots and resistance support transfer; the patent's 10-day root initiation is not transplant readiness. \textbf{[3,4]}}
\Step{8}{Transfer \& acclimate}{Move rooted material to a draining established-plant container. Start outdoor acclimation in semi-shade and increase sun gradually; protect against unsuitable cold. Use observation, not a fixed transplant day. \textbf{[3,4]}}
\end{minipage}\hfill\begin{minipage}[t]{.345\linewidth}\raggedright\vspace{0pt}
\Note{Gather}{Clean pruners, one small draining pot, a tray, and enough cactus medium or perlite to fill the pot.}
\Note{Milestones, not deadlines}{\begin{tabular}{@{}>{\raggedright\arraybackslash}p{.30\linewidth}>{\raggedright\arraybackslash}p{.60\linewidth}@{}}\textbf{Callus} & Several days; dry surface.\\\textbf{Roots} & 10 days at 24 C in the patent trial.\\\textbf{Transfer} & Roots plus resistance; no fixed date.\\\end{tabular}\par \textbf{[2,3,4]}}
\Note{Leaf option}{Use a whole leaf including attachment cells, with the attached end at the surface of slightly damp medium. General Sedum guidance reports roots and rosettes after several weeks; cultivar response is unverified. \textbf{[4]}}
\Note{Check 1}{Soft/dark base: check saturation, drainage and callusing; restart from sound tissue. \textit{Proposed.} \textbf{[3,4]}}
\Note{Check 2}{No resistance: stop tugging; check warmth and slight moisture, then allow time. \textit{Proposed.} \textbf{[2,3]}}
\Note{Check 3}{Scorch after transfer: return to semi-shade; resume gradual acclimation. \textit{Proposed.} \textbf{[3]}}
\end{minipage}\par\vspace{4pt}
\Panel{Ready for normal care}{Roots anchor the cutting and new growth continues. Return gradually to the established-plant care cues; propagation moisture is not a permanent watering schedule.\par Start date: \hrulefill\quad Method: \hrulefill\par First successful growth / conditions: \hrulefill}
\Sources{Evidence: 2 = SED-PAT USPTO cultivar patent; 3 = SED-MSU Montana State; 4 = SED-IA Iowa State. Proposed = proposed starting point, not a published ratio or household measurement. Full titles, URLs and scopes: entry sources.yaml and companion worksheets. Conversions rounded; source units authoritative.}{sedum-loves-fire / propagation}
